% Part: computability
% Chapter: recursive-functions
% Section: sequences

\documentclass[../../../include/open-logic-section]{subfiles}

\begin{document}

\olfileid[tr]{cmp}{rec}{seq}
\olsection{Diziler}

İlkel özyinelemeli fonksiyonlar kümesi dikkate değer ölçüde dayanıklıdır.
Fakat \emph{diziler} üzerinde çalışmak için yeterli bir yöntem geliştirdikten
sonra daha da fazlasını yapabileceğiz. Sonlu doğal sayı dizilerini doğal
sayılarla şu şekilde özdeşleştireceğiz: $\langle a_0,a_1,a_2,\dots,a_k\rangle$
dizisi
\[
  p_0^{a_0+1}\cdot p_1^{a_1+1}\cdot p_2^{a_2+1}\cdots p_k^{a_k+1}
\]
sayısına karşılık gelir. Örneğin $\langle2,7,3\rangle$ ile
$\langle2,7,3,0,0\rangle$ dizilerinin farklı sayısal kodları olmasını
güvenceye almak için üslere bir ekliyoruz. Hem $0$ hem $1$ boş diziyi
kodlayabilir; somutluk adına $\emptyseq$ simgesi $0$'ı göstersin.

Bu dizi kodlamasının çalışmasının nedeni Aritmetiğin Temel Teoremi denen
sonuçtur: Her $n\geq2$ doğal sayısı, $a_k\geq1$ olmak üzere
\[
  n=p_0^{a_0}\cdot p_1^{a_1}\cdots p_k^{a_k}
\]
biçiminde bir ve yalnızca bir yolla yazılabilir. Bu, eşlemenin
$\tuple{}(a_0,\dots,a_k)=\tuple{a_0,\dots,a_k}$ bire bir olduğunu güvenceye
alır: Farklı diziler farklı sayılara eşlenir ve her sayıya en çok bir dizi
karşılık gelir.

Şimdi bir dizinin uzunluğunu ve $i$. elemanını belirleme, diziye bir eleman
ekleme ve iki diziyi art arda birleştirme işlemlerinin tümünün ilkel
özyinelemeli olduğunu göstereceğiz.

\begin{prop}
  $s$ dizisinin uzunluğunu döndüren $\len{s}$ fonksiyonu ilkel
  özyinelemelidir.
\end{prop}

\begin{proof}
  $R(i,s)$ bağıntısı
  \[
    R(i,s) \text{ ancak ve ancak } p_i\mid s\land p_{i+1}\nmid s
  \]
  olarak tanımlansın. $R$ açıkça ilkel özyinelemelidir. $s$, boş olmayan
  bir dizinin kodu olduğunda, yani
  \[
    s=p_0^{a_0+1}\cdots p_k^{a_k+1}
  \]
  olduğunda $R(i,s)$ ancak ve ancak $p_i$, $s$'yi bölen en büyük asal,
  yani $i=k$ olduğunda geçerlidir. Dolayısıyla $s$ dizisinin uzunluğu,
  $p_i$ sayısının $s$'yi bölen en büyük asal olması ancak ve ancak $i+1$
  olur. Şöyle tanımlayabiliriz:
  \[
    \len{s}=\begin{cases}
      0 & \text{$s=0$ ya da $s=1$ ise}\\
      1+\bmin{i<s}{R(i,s)} & \text{aksi durumda.}
    \end{cases}
  \]
  Burada sınırlı en küçükleme kullanabiliriz; çünkü $s$ bir dizi koduyken
  $R(i,s)$ koşulunu sağlayan yalnızca bir $i$ vardır ve varsa $s$'den
  küçüktür.
\end{proof}

\begin{prop}
  $s$ dizisinin sonuna $a$ eklenmesiyle elde edilen diziyi döndüren
  $\fn{append}(s,a)$ fonksiyonu ilkel özyinelemelidir.
\end{prop}

\begin{proof}
  $\fn{append}$ şöyle tanımlanabilir:
  \[
    \fn{append}(s,a)=\begin{cases}
      2^{a+1} & \text{$s=0$ ya da $s=1$ ise}\\
      s\cdot p_{\len{s}}^{a+1} & \text{aksi durumda.}
    \end{cases}
  \]
\end{proof}

\begin{prop}
  $s$ dizisinin $i$. elemanını (ilk eleman $0$. diye adlandırılır),
  $i$ dizinin uzunluğundan büyük ya da ona eşitse $0$'ı döndüren
  $\fn{element}(s,i)$ fonksiyonu ilkel özyinelemelidir.
\end{prop}

\begin{proof}
  $a$, $s$ dizisinin $i$. elemanıdır ancak ve ancak $p_i^{a+1}$,
  $s$'yi bölen en büyük $p_i$ kuvvetidir; yani $p_i^{a+1}\mid s$ fakat
  $p_i^{a+2}\nmid s$ olur. Dolayısıyla
  \[
    \fn{element}(s,i)=\begin{cases}
      0 & \mbox{$i\geq\len{s}$ ise}\\
      \bmin{a<s}{(p_i^{a+2}\nmid s)} & \text{aksi durumda.}
    \end{cases}
  \]
\end{proof}

Yukarıda tanımlanan fonksiyonların resmî adları yerine daha kısa bir
gösterim kullanacağız. $\fn{element}(s,i)$ yerine $(s)_i$ yazacak;
$\tuple{s_0,\dots,s_k}$ ile
\[
  \fn{append}(\fn{append}(\dots\fn{append}(\emptyseq,s_0)\dots),s_k)
\]
ifadesini kısaltacağız. $s$ dizisinin uzunluğu $k$ ise elemanlarının
$(s)_0,\dots,(s)_{k-1}$ olduğuna dikkat edin.

\begin{prop}
İki diziyi art arda birleştiren $\fn{concat}(s,t)$ fonksiyonu ilkel
özyinelemelidir.
\end{prop}

\begin{proof}
  Aşağıdaki özelliği taşıyan bir $\fn{concat}$ fonksiyonu istiyoruz:
  \[
    \fn{concat}(\tuple{a_0,\dots,a_k},\tuple{b_0,\dots,b_l})=
    \tuple{a_0,\dots,a_k,b_0,\dots,b_l}.
  \]
  $t$ dizisinin ilk $n$ simgesini $s$ dizisinin sonuna ekleyen yardımcı bir
  $\fn{hconcat}(s,t,n)$ fonksiyonu kullanalım. Bu fonksiyon ilkel
  özyinelemeyle şöyle tanımlanır:
  \begin{align*}
    \fn{hconcat}(s,t,0) & = s\\
    \fn{hconcat}(s,t,n+1) & =
      \fn{append}(\fn{hconcat}(s,t,n),(t)_n)
    \intertext{Ardından $\fn{concat}$ fonksiyonunu şöyle tanımlayabiliriz:}
    \fn{concat}(s,t) & = \fn{hconcat}(s,t,\len{t}).
  \end{align*}
\end{proof}

$\fn{concat}(s,t)$ yerine $s\concat t$ yazacağız.

Bir dizinin sayısal kodunu, uzunluğu ve en büyük elemanı cinsinden
sınırlayabilmek yararlı olacaktır. $s$, uzunluğu $k$ olan ve her elemanı
bir $x$ sayısından küçük veya ona eşit olan bir dizi olsun. O zaman $s$'nin
en çok $k$ asal çarpanı vardır; bunların her biri en çok $p_{k-1}$ ve
$s$'nin asal çarpanlara ayrılışındaki üsleri en çok $x+1$'dir. Başka bir
deyişle
\[
  \fn{sequenceBound}(x,k)=p_{k-1}^{k\cdot(x+1)}
\]
diye tanımlarsak yukarıdaki $s$ dizisinin sayısal kodu en çok
$\fn{sequenceBound}(x,k)$ olur.

Diziler için böyle bir sınıra sahip olmak, sınırlı aramayla yeni fonksiyonlar
tanımlamamızı sağlar. Örneğin $\fn{concat}$ fonksiyonunu sınırlı aramayla
tanımlayabiliriz. Aradığımız nesnenin, yani birleştirilmiş dizinin sayısının
ilkel özyinelemeli bir \emph{belirtimini} ve ne kadar arayacağımıza ilişkin
bir sınırı yazmamız yeterlidir. Aşağıdaki tanım işe yarar:
\begin{align*}
  \fn{concat}(s,t) = {} & \bmin{v<\fn{sequenceBound}(s+t,\len{s}+\len{t})}{}\\
  & \quad(\len{v}=\len{s}+\len{t}\land{}\\
  & \qquad\bforall{i<\len{s}}{((v)_i=(s)_i)\land{}\\
  & \qquad\bforall{j<\len{t}}{((v)_{\len{s}+j}=(t)_j)})}.
\end{align*}

\begin{prob}
\[
  \fn{sconcat}(\tuple{s_0,\dots,s_k})=s_0\concat\dots\concat s_k
\]
özelliğini taşıyan ilkel özyinelemeli bir $\fn{sconcat}(s)$ fonksiyonu
bulunduğunu gösterin.
\end{prob}

\begin{prob}
\begin{align*}
  \fn{tail}(\emptyseq) & = 0 \text{ ve}\\
  \fn{tail}(\tuple{s_0,\dots,s_k}) & = \tuple{s_1,\dots,s_k}
\end{align*}
özelliklerini taşıyan ilkel özyinelemeli bir $\fn{tail}(s)$ fonksiyonu
bulunduğunu gösterin.
\end{prob}

\begin{prop}
  \ollabel{prop:subseq}
  $s$ dizisinin $i$. elemanından başlayan ve uzunluğu $n$ olan alt dizisini
  döndüren $\fn{subseq}(s,i,n)$ fonksiyonu ilkel özyinelemelidir.
\end{prop}

\begin{proof}
  Alıştırma.
\end{proof}

\begin{prob}
  \olref[cmp][rec][seq]{prop:subseq} önermesini ispatlayın.
\end{prob}

\end{document}

